\chapter{I backend CUDA}\label{cap:cuda}

L'implementazione su GPU è stata sviluppata attraverso tre kernel CUDA,
affiancati da un backend basato su cuBLAS, usato come riferimento esterno.
Le diverse versioni non sono state progettate come soluzioni indipendenti:
partono dalla stessa operazione locale, $Y = AX$, e introducono
progressivamente modifiche mirate alla mappatura del lavoro e agli accessi
alla memoria.

In particolare, il percorso seguito è il seguente:

\begin{itemize}
    \item \textbf{\kernelname{cuda\_naive}} costituisce la baseline CUDA.
    A ciascun thread viene assegnato un singolo elemento $Y_{ic}$:
    il thread fissa quindi una riga $i$ di $A$ e una colonna $c$ di $X$,
    e percorre l'intera dimensione interna $j$ per calcolarne il prodotto
    scalare. La mappatura è semplice e offre un primo termine di confronto,
    ma non sfrutta il riuso della riga di $A$ tra le $k$ colonne del
    multivettore.
\end{itemize}



\begin{itemize}
    \item \textbf{\kernelname{cuda\_warp}} cambia la granularità della
    parallelizzazione: non più un thread per elemento di $Y$, ma un
    \emph{warp} per un'intera riga di $Y$. Le 32 lane del warp si
    suddividono gli elementi della corrispondente riga di $A$ e ciascuna
    lane riutilizza il valore di $A$ appena caricato per aggiornare i
    $k$ risultati della riga. Al termine, i contributi prodotti dalle lane
    vengono sommati mediante primitive di warp shuffle.

    \begin{itemize}

    \item \textbf{\kernelname{cuda\_warp}} nella variante \code{xcol} mantiene
    esattamente la stessa organizzazione del calcolo, ma modifica la
    rappresentazione di $X$. Il blocco locale del multivettore viene
    convertito una sola volta, durante il preprocessing, da una disposizione
    in memoria per righe (\emph{row-major}) a una disposizione per colonne (\emph{column-major}).
    In questo modo quando le lane
    del warp avanzano su indici $j$ consecutivi, accedono anche a elementi
    consecutivi di $X$.

    \item \textbf{\kernelname{cuda\_warp\_smem}} affronta lo stesso problema degli
    accessi a $X$ senza però andarne a modificare il layout.
    Un blocco di thread carica
    \emph{cooperativamente} una porzione di $X$ dalla memoria globale alla
    \textbf{shared memory}. I warp appartenenti allo stesso blocco possono quindi
    riutilizzare quella stessa porzione durante il calcolo delle rispettive
    righe di $Y$.

    \end{itemize}
\end{itemize}

Il passaggio da \kernelname{cuda\_naive} a
\kernelname{cuda\_warp} modifica la mappatura del lavoro e introduce il
riuso di $A$ all'interno del warp. A questo punto è emerso un problema:
con $X$ nel normale layout row-major, le lane che avanzano lungo $j$ non
leggono elementi consecutivi del multivettore.

La variante \code{xcol} interviene sulla \emph{rappresentazione
del dato}, convertendo preventivamente $X$ in un layout più favorevole agli
accessi del warp, mentre \kernelname{cuda\_warp\_smem} conserva il layout
originario e interviene sul \emph{percorso del dato}, caricando e riutilizzando
$X$ attraverso la shared memory. Le due versioni permettono così di studiare
due strategie diverse per lo stesso collo di bottiglia, senza sovrapporle
nello stesso kernel.

Infine, si propone anche un'implementazione su GPU basata su cuBLAS, la libreria NVIDIA per algebra lineare ottimizzata per GPU:

\begin{itemize}
    \item \textbf{\kernelname{cublas}} non rappresenta un'ulteriore ottimizzazione
    del kernel sviluppato nel progetto. È un riferimento esterno, inserito
    dietro la stessa interfaccia degli altri backend, che permette di
    confrontare le implementazioni precedenti con una libreria GPU
    general-purpose sulla stessa macchina e nelle stesse condizioni di
    misura.
\end{itemize}

\section{Struttura comune}

Tutti i backend GPU rispettano lo stesso ciclo di vita a tre fasi introdotto
in \cref{sec:tre-fasi}, che permette di confrontare
le diverse implementazioni misurando, in ciascun caso, la stessa porzione
dell'esecuzione.

La distinzione fondamentale deriva dal diverso ruolo dei dati.Il blocco
locale di $A$ non cambia tra un'invocazione del prodotto e la successiva e
può quindi essere trasferito sulla GPU una sola volta. Al contrario,
la fetta locale di \(X\) deve essere resa disponibile a ciascun processo
a ogni invocazione del prodotto: le sue porzioni sono inizialmente collocate
sulla prima riga della griglia di processi e vengono replicate lungo ciascuna
colonna mediante MPI_Bcast., mentre il risultato parziale $Y$ deve essere riportato
sull'host prima della successiva \code{MPI\_Reduce}.
Di conseguenza, le operazioni comuni ai
backend CUDA sono distribuite nelle tre fasi nel modo seguente:

\begin{itemize}
    \item \textbf{\code{local\_gemm\_create}} prepara il backend prima dell'inizio
    delle ripetizioni cronometrate. In questa fase viene inizializzato il
    contesto CUDA, vengono allocati in VRAM i buffer destinati ad $A$, $X$
    e $Y$, vengono creati gli eventi CUDA usati per le misure e il blocco
    locale di $A$ viene copiato dall'host al device. Poiché $A$ rimane
    invariata per tutte le invocazioni successive, questo trasferimento
    viene eseguito una sola volta. L'intera fase appartiene al
    preprocessing e non entra nel tempo ufficiale del prodotto. I suoi
    costi principali vengono comunque registrati separatamente come
    inizializzazione del device, allocazione della memoria e trasferimento
    H2D di $A$.

    \item \textbf{\code{local\_gemm}} rappresenta invece il cammino percorso a ogni
    invocazione. Prima del calcolo, il nuovo blocco di $X$ viene trasferito
    dall'host alla GPU. Segue l'esecuzione del kernel CUDA, oppure della
    chiamata cuBLAS nel backend di riferimento. Infine il risultato
    parziale $Y$ viene copiato nuovamente verso l'host. Il tempo
    \code{t\_local} comprende quindi l'intera sequenza
    H2D di $X$, calcolo e D2H di $Y$.

    All'interno della stessa funzione, tre coppie distinte di
    \code{cudaEvent} delimitano separatamente il trasferimento H2D di $X$,
    il solo calcolo sul device e il trasferimento D2H di $Y$. Si ottengono
    in questo modo tre misurazioni: rispettivamente $t_{\mathrm{H2D}\,X}$,
    $t_{\mathrm{kernel}}$ e $t_{\mathrm{D2H}\,Y}$.
    Questa separazione è necessaria perché la specifica del progetto
    richiede di escludere dal tempo ufficiale i trasferimenti da e verso
    la scheda: per i backend CUDA è quindi $t_{\mathrm{kernel}}$, e non
    l'intero $t_{\mathrm{local}}$, a rappresentare il contributo del
    calcolo locale alla misura prestazionale.

    \item \textbf{\code{local\_gemm\_destroy}} conclude infine il ciclo di vita del
    backend, liberando i buffer allocati sulla GPU, gli eventi CUDA e il
    contesto mantenuto dall'implementazione. Anche questa fase avviene al
    di fuori delle ripetizioni misurate.
\end{itemize}

La stessa suddivisione viene mantenuta in tutti i backend GPU. Cambia quindi
il modo in cui viene eseguito il prodotto all'interno della seconda fase,
ma non cambiano né la posizione dei trasferimenti né i confini delle misure.
Il confronto tra i kernel può così essere effettuato sul solo
$t_{\mathrm{kernel}}$, mentre $t_{\mathrm{local}}$ rimane disponibile per
osservare anche il costo reale dell'intera invocazione lato GPU.

\section{Misurazione del backend CUDA}
\label{sec:misurazione-cuda}

La natura asincrona dell'esecuzione CUDA richiede di distinguere il tempo
osservato dall'host dal tempo trascorso effettivamente sul device. Il lancio
di un kernel, infatti, restituisce il controllo alla CPU senza attendere
necessariamente il completamento dell'operazione. Un cronometro host posto
semplicemente prima e dopo il lancio misurerebbe semplicemente il costo
della chiamata al runtime CUDA, non il tempo di esecuzione del kernel,
ciò che invece ci interessa.

Per isolare le singole operazioni eseguite sul device, i backend CUDA
utilizzano coppie di \textbf{\code{cudaEvent}} registrate sullo stesso stream delle
operazioni misurate: in questo modo gli eventi permettono di determinare
separatamente il tempo necessario al trasferimento H2D di $X$,
all'esecuzione del calcolo e al trasferimento D2H di $Y$.

Il tempo della fase locale osservato dal driver, indicato con
$t_{\mathrm{local}}$, ha invece un significato diverso: misura l'intera
invocazione del backend e comprende quindi sia il lavoro eseguito sul device
sia l'overhead necessario per richiederlo e sincronizzarlo.

Schematizzando il calcolo, otteniamo:

\begin{equation}
t_{\mathrm{local}}
=
t_{\mathrm{H2D}\,X}
+
t_{\mathrm{kernel}}
+
t_{\mathrm{D2H}\,Y}
+
t_{\mathrm{overhead}},
\label{eq:scomposizione-locale}
\end{equation}

dove $t_{\mathrm{overhead}}$ non viene misurato direttamente, ma rappresenta
il residuo fra il tempo complessivo osservato dall'host e i tre intervalli
misurati mediante eventi CUDA. Esso può comprendere il costo delle chiamate
al runtime, delle sincronizzazioni e della gestione della strumentazione.

La distinzione è rilevante anche per la valutazione delle prestazioni,
perché come richiesto dalla specifica del progetto, i trasferimenti tra host e
device non devono contribuire al tempo utilizzato per calcolare i FLOPS.
Per i backend CUDA il contributo computazionale è quindi
$t_{\mathrm{kernel}}$, mentre $t_{\mathrm{H2D}\,X}$ e
$t_{\mathrm{D2H}\,Y}$ vengono conservati come misure separate e analizzati
successivamente.

Oltre al tempo, il backend comunica al driver anche il numero di byte
trasferiti nelle tre operazioni principali: la copia iniziale di $A$, la
copia di $X$ eseguita a ogni invocazione e la copia di ritorno di $Y$.
Tempo e volume trasferito permettono di ricavare una \textbf{banda effettiva}, calcolata come:

\begin{equation}
B_{\mathrm{eff}} = \frac{D}{t},
\end{equation}

dove $D$ è il numero di byte trasferiti e $t$ il tempo della corrispondente
copia.

Il solo tempo di trasferimento sarebbe infatti poco informativo: una copia
più grande richiede naturalmente più tempo di una più piccola. La banda
effettiva permette invece di confrontare trasferimenti di dimensioni
differenti e di valutare quanto efficacemente venga utilizzato il collegamento
tra host e GPU.

\section{Preprocessing e occupazione della GPU}
\label{sec:preprocessing-cuda}

Anche il costo di creazione del backend viene registrato, pur rimanendo
esterno alla regione utilizzata per la misura ufficiale. Il tempo
$t_{\mathrm{setup}}$ viene distinto in tre contributi, perché corrispondono
a operazioni di natura differente.

Il primo è l'inizializzazione dell'ambiente CUDA, necessaria affinché il
processo possa utilizzare il device. Si tratta principalmente di un costo
di avvio e non dipende direttamente dalle dimensioni del prodotto da
calcolare.

Il secondo contributo riguarda l'allocazione in memoria device dei buffer
necessari per $A$, $X$ e $Y$. Questo costo è legato alla preparazione delle
risorse del backend e non rappresenta un trasferimento di dati.

Il terzo contributo è invece la copia H2D del blocco locale di $A$.
Poiché $A$ rimane invariata durante tutte le successive invocazioni del
prodotto, questa copia viene effettuata una sola volta in
\code{local\_gemm\_create}. Il suo costo viene quindi escluso dal cammino
cronometrato, ma misurato separatamente.

Questa distinzione permette di quantificare direttamente il vantaggio della
persistenza di $A$ sul device: se il suo trasferimento fosse ripetuto a ogni
chiamata di \code{local\_gemm}, il relativo costo entrerebbe nel percorso
critico insieme alla copia di $X$ e a quella di $Y$.

\subsection{Gestione della memoria disponibile}

La disponibilità della memoria device non viene verificata mediante una
stima preventiva utilizzata per decidere se procedere con il calcolo.
Il backend tenta invece direttamente le allocazioni necessarie mediante
\code{cudaMalloc} e, a questo punto, se si verifica un errore di allocazione, viene
registrato un messaggio di errore e si procede con il
fallimento della creazione del backend.

Questa scelta evita di basare la decisione su una precedente interrogazione
della memoria libera: infatti, in presenza di più processi MPI che condividono la
stessa GPU, il valore restituito da una precedente osservazione
dello stato della memoria può cambiare prima dell'allocazione effettiva.
In caso di errore, \code{cudaMemGetInfo} viene comunque utilizzata per
produrre una diagnostica che riporta la memoria richiesta e quella
disponibile.

I principali buffer persistenti richiedono complessivamente una quantità di byte pari a:

\begin{equation}
\left(
    \mloc \cdot \code{lda}
    +
    \nloc \cdot \code{ldx}
    +
    \mloc \cdot \code{ldy}
\right)
\cdot \code{sizeof(scalar\_t)}
\end{equation}

Poiché nel problema considerato $k$ è piccolo rispetto alle dimensioni
della matrice, il termine associato ad $A$ è normalmente dominante.

Aumentare il numero di processi MPI e scegliere opportunamente la griglia
bidimensionale riduce le dimensioni dei blocchi locali e, di conseguenza,
la quantità di memoria device richiesta da ciascun processo.
\section{\kernelname{cuda\_naive}}\label{sec:naive}

La prima implementazione CUDA adotta la mappatura più diretta possibile:
a ciascun thread viene assegnato un singolo elemento del risultato locale
$Y$.

Indicando con $t$ l'identificativo globale del thread, gli indici
dell'elemento da calcolare sono

\begin{equation}
    i = \left\lfloor \frac{t}{k} \right\rfloor,
    \qquad
    c = t \bmod k.
\end{equation}

Il thread $t$ è quindi responsabile di $Y_{ic}$ e ne calcola il valore
percorrendo interamente la dimensione interna del prodotto:

\begin{equation}
    Y_{ic}
    =
    \sum_{j=0}^{\nloc-1}
    A_{ij}X_{jc}.
\end{equation}

Ogni thread mantiene un unico accumulatore e non comunica con gli altri
thread.

Il kernel costituisce per questo una baseline semplice, utile sia
per verificare il corretto funzionamento dell'intera pipeline CUDA, sia
come riferimento interno rispetto al quale valutare le ottimizzazioni
successive.

La semplicità della mappatura comporta però una prima inefficienza
fondamentale. Infatti, per una stessa riga $i$, gli elementi:

\[
Y_{i0},Y_{i1},\ldots,Y_{i,k-1}
\]

vengono calcolati da thread distinti. Durante una stessa iterazione $j$,
tutti questi thread hanno bisogno dello stesso coefficiente $A_{ij}$,
ma ciascuno lo utilizza soltanto per la propria colonna di $Y$.
Il kernel non esprime quindi alcun riuso di $A_{ij}$ all'interno del
thread: lo stesso elemento di $A$ compare logicamente in $k$ operazioni
di caricamento differenti.

Questo è il principale limite che motiva il backend successivo: in
\kernelname{cuda\_warp}, un valore di $A$ viene invece caricato una volta
da una lane e riutilizzato per aggiornare tutti i $k$ risultati associati
alla stessa riga.

La mappatura del kernel assegna thread consecutivi alle colonne consecutive di
$Y$. Per un $j$ fissato, tutte le lane che lavorano sulla stessa riga $i$ di
$Y$ richiedono quindi lo stesso elemento $A_{ij}$. Le lane associate alla
riga successiva richiedono invece $A_{i+1,j}$, che in memoria row-major si
trova a distanza \code{lda} elementi da $A_{ij}$.

Per $k<32$, cioè quando una riga di $Y$ contiene meno elementi della
dimensione di un warp, un singolo warp comprende lane associate a più righe
di $Y$ e coinvolge quindi più gruppi di accessi ad $A$ distanziati tra loro.

Gli accessi a $X$ hanno invece un comportamento differente. Per ogni gruppo
di lane associato alla stessa riga di $Y$, thread consecutivi leggono gli
elementi consecutivi

\[
X_{j0}, X_{j1}, \ldots, X_{j,k-1},
\]

che sono contigui in memoria nel layout row-major. Se il warp comprende più
righe di $Y$, questa stessa sequenza di accessi a $X$ si ripete.

All'aumentare di $k$, un numero crescente di lane del warp lavora sulla
stessa riga di $Y$, fino al caso $k=32$, in cui un intero warp corrisponde
esattamente a una riga. A ogni iterazione $j$, tutte le lane richiedono
quindi lo stesso elemento $A_{ij}$ e, contemporaneamente, leggono i 32
elementi consecutivi

\[
X_{j0}, \ldots, X_{j31}.
\]

Il parametro $k$ influenza quindi direttamente il pattern di accesso del
kernel. Questo aspetto sarà utile nel confronto con
\kernelname{cuda\_warp}, dove la mappatura viene rovesciata: un intero
warp viene assegnato a una sola riga di $Y$ e le lane vengono distribuite
lungo la dimensione $j$.


\section{\kernelname{cuda\_warp}: un warp per riga}
\label{sec:warp}

Il secondo backend modifica radicalmente l'approccio alla
parallelizzazione. In \kernelname{cuda\_naive} ogni thread era responsabile
di un singolo elemento $Y_{ic}$, ma in \kernelname{cuda\_warp}, invece, un
intero warp collabora ed è responsabile del calcolo di una riga di $Y$.

Fissata la riga $i$, le 32 lane del warp si dividono la dimensione interna
$j$. La lane $l$ visita gli indici:

\[
j = l,\; l+32,\; l+64,\ldots
\]

fino a raggiungere $\nloc$.

A ogni iterazione, lane consecutive accedono
quindi a elementi consecutivi della stessa riga della matrice $A$:

\[
A_{i,j}, A_{i,j+1}, \ldots, A_{i,j+31}.
\]

La mappatura produce così un pattern di accesso contiguo, favorevole alla
coalescenza delle letture dalla memoria globale. A differenza di quanto
accadeva in \kernelname{cuda\_naive}, questo comportamento non dipende da
$k$, perché è la dimensione $j$, e non quella delle colonne del
multivettore, a essere distribuita fra le lane.

La nuova mappatura consente soprattutto di sfruttare esplicitamente il
riuso degli elementi di $A$. Una volta caricato un elemento:

\[
a = A_{ij},
\]

la lane non lo utilizza per una sola colonna del risultato, ma per
aggiornare tutti i $k$ accumulatori associati alla riga di $Y$:

\[
\begin{aligned}
acc[0]   &\mathrel{+}= aX_{j0},\\
acc[1]   &\mathrel{+}= aX_{j1},\\
          &\vdots\\
acc[k-1] &\mathrel{+}= aX_{j,k-1}.
\end{aligned}
\]

Un singolo caricamento di $A_{ij}$ viene quindi riutilizzato nelle $k$
moltiplicazioni che coinvolgono lo stesso elemento della matrice, ed è questa la principale
 differenza rispetto al primo approccio.

Al termine del ciclo su $j$, ciascuna lane possiede però soltanto una parte
del risultato. Per ogni colonna $c$, la lane $l$ ha accumulato i termini
corrispondenti agli indici:

\[
j = l,\;l+32,\;l+64,\ldots
\]

e il valore completo può essere scritto come

\begin{equation}
Y_{ic}
=
\sum_{l=0}^{31} acc_l[c].
\end{equation}

I 32 contributi vengono quindi ridotti all'interno del warp mediante
\code{\_\_shfl\_down\_sync}. Al termine, la lane $0$
possiede i $k$ valori completi della riga e li scrive in $Y$.

\subsection{Gestione di $k$}

I valori di $k$ richiesti dalla traccia, $3$, $6$, $8$, $20$ e $32$, sono
implementati mediante cinque istanze template distinte del kernel. In
queste versioni $k$ è quindi noto a compile time: la dimensione
dell'array degli accumulatori è costante e i cicli sulle colonne possono
essere srotolati dal compilatore.

La stessa interfaccia supporta comunque valori arbitrari di $k$. Quando
$k$ non appartiene all'insieme precedente viene utilizzato un kernel
runtime che elabora quattro colonne alla volta
(\code{RUNTIME\_TILE = 4}). In questo modo il numero di accumulatori
simultaneamente mantenuti da ciascuna lane rimane limitato, pur continuando
a riutilizzare ogni elemento di $A$ per più colonne.

\subsection{La maschera dello shuffle}

Un dettaglio di correttezza riguarda la maschera utilizzata nelle
operazioni di shuffle. Il kernel usa la maschera completa
\code{0xffffffff}, perché ogni warp valido porta tutte le 32 lane fino
alla riduzione. Il controllo sul limite delle righe viene infatti eseguito
uniformemente sull'intero warp: se la riga assegnata non esiste, tutte le
lane terminano prima; altrimenti tutte raggiungono la
\code{\_\_shfl\_down\_sync}.

Questo rimane vero anche quando $\nloc < 32$: le lane il cui indice non
corrisponde ad alcun $j$ mantengono semplicemente accumulatori nulli, ma
partecipano comunque alla riduzione.

\subsection{Il problema residuo: gli accessi a $X$}\label{sec:difetto-x}

La distribuzione di $j$ fra le lane rende favorevoli gli accessi ad $A$,
ma determina contemporaneamente un pattern meno favorevole per $X$ quando
il multivettore è memorizzato in row-major.

Consideriamo un'iterazione nella quale le 32 lane del warp stanno
elaborando gli indici consecutivi

\[
j_0,\;j_0+1,\ldots,j_0+31.
\]

Durante l'aggiornamento di una stessa colonna $c$, le lane leggono quindi

\[
X_{j_0,c},\;
X_{j_0+1,c},\;
\ldots,\;
X_{j_0+31,c}.
\]

Si tratta di elementi appartenenti alla stessa colonna di $X$, ma a 32
righe differenti. Poiché $X$ è memorizzata in row-major, due elementi
consecutivi di questa sequenza non sono adiacenti in memoria: i loro
indirizzi distano \code{ldx} elementi, cioè una quantità di byte pari a:

\[
\code{ldx}\cdot \code{sizeof(scalar\_t)}
\]

dove ricordiamo che nel caso compatto, utilizzato dal progetto, si ha sempre $\code{ldx}=k$.

Il comportamento è quindi opposto a quello delle letture di $A$:
per la matrice $A$, lane consecutive leggono elementi consecutivi della stessa riga,
mentre per $X$ durante il calcolo di una colonna $c$, lane consecutive leggono
elementi della stessa colonna ma appartenenti a righe differenti.

All'aumentare di $k$ aumenta anche la distanza in memoria fra gli elementi
di $X$ richiesti da lane consecutive. Le richieste del warp risultano
quindi progressivamente più sparse e possono richiedere un numero maggiore
di transazioni di memoria rispetto a un accesso contiguo.

Il kernel ha quindi risolto il principale limite di
\kernelname{cuda\_naive}, cioè il mancato riuso esplicito di $A$, ma ha
lasciato un nuovo collo di bottiglia nel pattern di accesso a $X$.
Le due varianti introdotte nelle sezioni successive affrontano esattamente
questo problema senza modificare la mappatura warp-per-row.

\section{Primo rimedio: cambio di layout di $X$}
\label{sec:x-layout}

Il problema descritto nella sezione precedente nasce dall'incompatibilità
fra la mappatura delle lane e il layout row-major di $X$. Le lane di
\kernelname{cuda\_warp} avanzano lungo la dimensione $j$, mentre nel layout
row-major sono gli elementi che differiscono per l'indice di colonna $c$
a essere contigui in memoria.

Il primo rimedio consiste quindi nel modificare la rappresentazione di $X$
senza cambiare la mappatura del kernel. Durante il preprocessing, il blocco
locale del multivettore viene riscritto per colonne, in modo che

\begin{equation}
X_{jc}
=
\code{X\_loc}\bigl[c \cdot \nloc + j\bigr].
\label{eq:xcol}
\end{equation}

In questa rappresentazione, tutti gli elementi appartenenti a una stessa
colonna di $X$ sono consecutivi in memoria. Di conseguenza, quando le 32
lane del warp elaborano gli indici

\[
j_0,\;j_0+1,\ldots,j_0+31
\]

per una stessa colonna $c$, leggono anche 32 posizioni consecutive del
buffer:

\[
X_{j_0,c},\;
X_{j_0+1,c},\;
\ldots,\;
X_{j_0+31,c}.
\]

Lo stride fra gli indirizzi richiesti da lane consecutive passa quindi da
$k$ elementi (nel layout row-major) a un singolo elemento (in questo layout).
La mappatura warp-per-row e il riuso di $A$ introdotti da \kernelname{cuda\_warp}
rimangono invariati: cambia esclusivamente il modo in cui $X$ è organizzata
e indirizzata in memoria.






\section{\kernelname{cuda\_warp\_smem}: staging di $X$ in shared memory}
\label{sec:smem}

Il secondo rimedio conserva sia la mappatura warp-per-row di
\kernelname{cuda\_warp} sia il normale layout row-major di $X$.
La differenza riguarda il percorso seguito dagli elementi del multivettore:
invece di essere letti direttamente dalla global memory durante il calcolo,
vengono prima caricati cooperativamente in \textbf{shared memory} e successivamente
consumati dai warp del blocco.

Il percorso di un elemento di $X$ diventa quindi:

\[
\text{global memory}
\;\longrightarrow\;
\text{shared memory}
\;\longrightarrow\;
\text{registri},
\]

mentre in \kernelname{cuda\_warp} il passaggio intermedio non era presente.

La struttura del calcolo rimane per il resto invariata. Ogni warp è ancora
responsabile di una riga di $Y$, le sue 32 lane si dividono gli indici $j$,
ogni elemento di $A$ viene riutilizzato per aggiornare i $k$ accumulatori
e la riduzione finale avviene mediante \code{\_\_shfl\_down\_sync}.

\subsection{Caricamento cooperativo di $X$}

Il caricamento del tile di $X$ non utilizza la stessa mappatura con cui
viene successivamente eseguito il prodotto. In questa prima fase
partecipano tutti i thread del blocco e gli elementi del tile vengono
assegnati mediante un indice lineare.

Indicando con $q$ tale indice, vengono ricavati

\[
r = \left\lfloor \frac{q}{k} \right\rfloor,
\qquad
c = q \bmod k,
\]

e il thread carica l'elemento $X_{rc}$ nella corrispondente posizione del
tile in shared memory.

Con $X$ compatta in row-major, thread consecutivi ricevono valori
consecutivi di $q$ e leggono quindi posizioni consecutive del buffer.
Il caricamento dalla global memory segue così l'ordine naturale con cui
$X$ è memorizzata, evitando il pattern strided che caratterizzava
\kernelname{cuda\_warp}.

Una volta completato il caricamento cooperativo, una
\code{\_\_syncthreads()} garantisce che il tile sia interamente disponibile
prima che i warp inizino a utilizzarlo. La shared memory è infatti una
risorsa comune a tutti i thread dello stesso blocco.

A questo punto il calcolo riprende la stessa struttura di
\kernelname{cuda\_warp}. Per ogni warp, la lane $l$ elabora gli indici

\[
j=l,\;l+32,\;l+64,\ldots
\]

ma gli elementi di $X$ corrispondenti vengono ora letti dal tile in shared
memory anziché direttamente dalla global memory.

\subsection{Riuso del tile fra i warp del blocco}

L'introduzione della shared memory produce inoltre un secondo effetto.
I warp appartenenti allo stesso blocco calcolano righe differenti di $Y$,
e utilizzano quindi righe differenti di $A$, ma tutti devono attraversare
la stessa matrice $X$.

Nel backend \kernelname{cuda\_warp}, ciascun warp emette le proprie letture
di $X$. In \kernelname{cuda\_warp\_smem}, invece, un tile viene caricato una
sola volta dall'intero blocco e viene poi utilizzato da tutti i suoi warp.

Se $W$ indica il numero di warp per blocco, il numero di elementi di $X$
caricati esplicitamente dalla global memory passa quindi, idealmente, da

\[
m_{\mathrm{loc}}\,n_{\mathrm{loc}}\,k
\]

a

\[
\left\lceil\frac{m_{\mathrm{loc}}}{W}\right\rceil
n_{\mathrm{loc}}\,k.
\]

Per esempio, con il valore predefinito \code{BLOCK\_THREADS=256}, il blocco contiene
$W=8$ warp.

Questa riduzione descrive il numero di letture esplicitamente richieste
dal kernel alla gerarchia di memoria globale. Non implica però che il
traffico effettivo verso la DRAM si riduca dello stesso fattore, perché
una parte delle letture di \kernelname{cuda\_warp} può già essere servita
dalle cache della GPU. Il vantaggio effettivo della shared memory deve
quindi essere determinato sperimentalmente.

\subsection{Bank conflict e padding della shared memory}
\label{sec:bank-conflict}

Una volta caricato il tile, il problema non è più la coalescenza degli
accessi alla global memory, ma il modo in cui gli indirizzi richiesti dalle
lane vengono distribuiti sui banchi della shared memory.

La shared memory è organizzata in 32 banchi, e parole consecutive da
32 bit vengono assegnate a banchi consecutivi. Se più lane richiedono
indirizzi differenti appartenenti allo stesso banco, gli accessi non
possono essere serviti simultaneamente e vengono parzialmente
serializzati.

Durante il caricamento cooperativo del tile, thread consecutivi scrivono
elementi consecutivi della shared memory. Questo pattern segue
linearmente il buffer e non introduce il problema che interessa il
successivo consumo del tile.

Durante il calcolo, invece, la lane $l$ accede alla posizione

\[
X_{\mathrm{tile}}
\bigl[l\cdot\code{tile\_row\_stride}+c\bigr],
\]

dove

\[
\code{tile\_row\_stride}
=
k+\code{SMEM\_PAD}.
\]

Lane consecutive leggono quindi elementi separati da un'intera riga del
tile. Se questo stride produce una mappatura periodica sugli stessi banchi,
si genera un bank conflict.

Il caso $k=32$ in doppia precisione è particolarmente sfavorevole senza
padding. Con

\[
\code{tile\_row\_stride}=32,
\]

la distanza fra gli elementi letti da lane consecutive corrisponde a
$64$ parole da 32 bit e quindi

\[
64 \equiv 0 \pmod{32}.
\]

Gli accessi ricadono così ripetutamente sugli stessi banchi.

Aggiungendo un elemento di padding per riga si ottiene invece

\[
\code{tile\_row\_stride}=33,
\]

e la distanza diventa

\[
66 \equiv 2 \pmod{32},
\]

distribuendo gli accessi sui diversi banchi e rimuovendo il caso
patologico precedente.

Il padding non rappresenta tuttavia una scelta universalmente ottimale
per ogni $k$. Il suo effetto dipende dallo stride risultante. Per questo
\code{SMEM\_PAD} è un parametro di compilazione e il progetto mantiene
separate le configurazioni con e senza padding, così da misurarne
direttamente l'effetto.

\subsection{Scelta della dimensione del tile}
\label{sec:piano-tile}

Il blocco locale di $X$ non viene trasferito interamente in shared memory.
La dimensione della shared memory disponibile per blocco è limitata e
contribuisce, insieme ai registri e al numero di thread, a determinare
quanti blocchi possono risiedere contemporaneamente su uno SM.

Il backend divide quindi $X$ in tile contenenti
\code{x\_rows\_per\_tile} righe. Indicando questa quantità con $R$, la
memoria richiesta da un tile è

\[
R\,
\bigl(k+\code{SMEM\_PAD}\bigr)\,
\code{sizeof(scalar\_t)}.
\]

Il valore di $R$ modifica la quantità di shared memory occupata da ogni
blocco, ma non il volume complessivo di $X$ che quel blocco deve
attraversare. Un tile più piccolo produce semplicemente un numero maggiore
di iterazioni del ciclo di tiling.

La dimensione del tile influenza invece direttamente la frequenza delle
sincronizzazioni. Ogni tile richiede due chiamate a
\code{\_\_syncthreads()}: la prima impedisce che il tile precedente venga
sovrascritto mentre qualche warp lo sta ancora utilizzando, la seconda
garantisce che il nuovo tile sia stato completamente caricato prima di
iniziare il calcolo. Aumentando il numero di righe per tile diminuisce
quindi il numero di barriere incontrate durante l'attraversamento di $X$.

Non è tuttavia possibile aumentare arbitrariamente il tile. Una maggiore
quantità di shared memory per blocco può ridurre il numero di blocchi
residenti contemporaneamente su uno SM e quindi l'occupancy. La scelta
implementata nel backend cerca perciò un compromesso: utilizzare un tile
sufficientemente grande senza fare della shared memory un ulteriore
vincolo alla residenza dei blocchi.

La funzione \code{plan\_tile} esegue questa scelta una sola volta durante
\code{local\_gemm\_create}, quindi fuori dalla regione cronometrata.

Come primo passo viene interrogata
\code{cudaOccupancyMaxActiveBlocksPerMultiprocessor} assumendo nulla la
shared memory dinamica. Il risultato rappresenta il numero di blocchi per
SM consentito dal kernel compilato in assenza del vincolo aggiuntivo
introdotto dal tile.

La shared memory disponibile per SM viene quindi suddivisa rispetto a
questo numero di blocchi, ottenendo un budget iniziale per ciascun blocco.
Da tale budget viene ricavato il numero di righe che possono essere
contenute nel tile.

Il valore così ottenuto viene infine verificato interrogando nuovamente
l'API di occupancy con la quantità effettiva di shared memory richiesta.
Se il numero di blocchi residenti risulta inferiore all'obiettivo iniziale,
la dimensione del tile viene progressivamente ridotta.

La pianificazione viene eseguita sull'istanza concreta del kernel
specializzato per il valore di $k$ corrente. Questa distinzione è
necessaria perché le diverse istanze possono avere richieste differenti
di registri e quindi differenti limiti di occupancy.

Il parametro \code{TILE\_GRANULARITY} controlla infine
l'arrotondamento del numero di righe scelto dal piano. Il valore
predefinito è $32$, pari alla dimensione del warp.

Un tile con un numero di righe multiplo di $32$ permette a tutte le lane
di partecipare a ciascun passo completo del ciclo su $j$. Questo non è
però un requisito di correttezza: il kernel gestisce anche tile con un
numero di righe non multiplo di $32$.

L'arrotondamento introduce quindi un compromesso. Una granularità elevata
può lasciare inutilizzata una parte del budget di shared memory e
aumentare il numero di tile, mentre una granularità più fine può produrre
un ultimo passo nel quale alcune lane non hanno elementi da elaborare.
Il parametro viene mantenuto configurabile proprio per poter misurare
questo compromesso.

\subsection{Sincronizzazione e warp inattivi}

L'uso di \code{\_\_syncthreads()} introduce un vincolo di correttezza che
non era presente in \kernelname{cuda\_warp}. La barriera coinvolge tutti
i thread del blocco e può essere superata soltanto quando tutti i thread
che partecipano al blocco l'hanno raggiunta.

L'ultimo blocco della griglia può contenere più warp di quanti siano
necessari per coprire le ultime righe di $Y$. In
\kernelname{cuda\_warp} questi warp possono terminare immediatamente
quando la riga assegnata supera $\mloc$.

In \kernelname{cuda\_warp\_smem} questo non è possibile prima del
completamento del ciclo sui tile. Se alcuni warp uscissero mentre gli
altri continuano a eseguire le \code{\_\_syncthreads()}, i thread
rimanenti potrebbero attendere indefinitamente a una barriera che gli
altri non raggiungeranno mai.

Per questo il kernel associa a ciascun warp un flag \code{active}.
Anche i warp ai quali non corrisponde una riga valida partecipano al
caricamento cooperativo e attraversano tutte le sincronizzazioni del
blocco. Soltanto l'accumulo dei prodotti e la successiva scrittura di
$Y$ vengono eseguiti dai warp attivi.

Terminato l'ultimo tile, quando non rimangono più barriere di blocco,
i warp inattivi possono infine uscire prima della riduzione finale.

\section{Due rimedi allo stesso problema}
\label{sec:due-rimedi}

Le soluzioni descritte nelle \cref{sec:x-layout,sec:smem} nascono dallo
stesso problema: in \kernelname{cuda\_warp}, con $X$ row-major, lane
consecutive leggono elementi della stessa colonna di $X$ separati da uno
stride pari a $k$ elementi.

I due approcci intervengono però in punti differenti.

La variante \code{X\_LAYOUT=column} modifica la rappresentazione del dato
prima dell'esecuzione del kernel. Disponendo consecutivamente gli elementi
di una stessa colonna di $X$, la mappatura warp-per-row può essere mantenuta
inalterata e le lane trovano direttamente in global memory gli elementi
richiesti in posizioni consecutive. Il costo associato a questa soluzione
è la conversione preliminare del multivettore.

\kernelname{cuda\_warp\_smem}, invece, conserva $X$ nel normale layout
row-major e modifica il percorso attraverso il quale i dati raggiungono le
lane. Gli elementi vengono prima caricati cooperativamente e in modo
contiguo dalla global memory alla shared memory, quindi letti dal tile
durante il calcolo. Questa soluzione evita la conversione del dato, ma
introduce il costo dello staging, delle sincronizzazioni e dell'utilizzo
della shared memory.

Le due implementazioni rappresentano quindi due strategie alternative
studiate separatamente nel progetto: la prima modifica il layout del dato,
la seconda modifica il meccanismo di accesso mantenendone invariata la
rappresentazione.

La combinazione dei due approcci non è stata inclusa tra i backend
sperimentali, poiché l'obiettivo del confronto è isolare separatamente
l'effetto delle due strategie rispetto alla stessa baseline
\kernelname{cuda\_warp}.

\section{\kernelname{cublas}: il riferimento esterno}
\label{sec:cublas}

L'ultimo backend GPU non introduce un nuovo kernel sviluppato nel progetto,
ma utilizza la routine GEMM della libreria cuBLAS. Il suo ruolo è fornire un
riferimento esterno con cui confrontare le implementazioni CUDA precedenti.

Perché il confronto sia significativo, cuBLAS viene inserita dietro la stessa
interfaccia degli altri backend. La creazione del contesto, le allocazioni
device e il trasferimento iniziale di $A$ avvengono quindi in
\code{local\_gemm\_create}, mentre ogni invocazione di \code{local\_gemm}
esegue, nello stesso ordine degli altri backend GPU, il trasferimento H2D di
$X$, il prodotto sul device e il trasferimento D2H di $Y$.

Anche la misura segue gli stessi confini: il tempo del solo prodotto cuBLAS
viene delimitato mediante \code{cudaEvent}, mentre i trasferimenti vengono
misurati separatamente. Il confronto con i kernel sviluppati nel progetto
avviene quindi nelle stesse condizioni sperimentali.

\subsection{Adattamento fra row-major e column-major}

Il progetto memorizza le matrici in row-major, mentre l'interfaccia GEMM
tradizionale di cuBLAS interpreta i propri operandi in column-major.
Non è tuttavia necessario convertire o trasporre fisicamente i dati.

Cioè gli stessi dati che rappresentano una matrice in row-major, se interpretati
come column-major senza modificarne il contenuto in memoria, rappresentano
la sua trasposta.
Cambia soltanto la convenzione con cui gli indici vengono associati agli indirizzi.

Di conseguenza, i buffer del progetto possono essere reinterpretati da
cuBLAS come

\[
A_{\mathrm{cm}} = A_{\mathrm{rm}}^T,
\qquad
X_{\mathrm{cm}} = X_{\mathrm{rm}}^T,
\qquad
Y_{\mathrm{cm}} = Y_{\mathrm{rm}}^T.
\]


Applicando questa reinterpretazione al prodotto desiderato si ottiene:

\[
Y_{\mathrm{cm}}
=
Y_{\mathrm{rm}}^{T}
=
\left(A_{\mathrm{rm}}X_{\mathrm{rm}}\right)^{T}
=
X_{\mathrm{rm}}^{T}A_{\mathrm{rm}}^{T}
=
X_{\mathrm{cm}}A_{\mathrm{cm}}.
\]

Quindi la chiamata cuBLAS deve utilizzare $X$ come primo operando e $A$ come
secondo. Nella convenzione utilizzata internamento da cuBLAS si ha:

\[
C_{M\times N}=A_{M\times K}B_{K\times N},
\]

con:

\[
C_{M\times N}=Y_{\mathrm{cm}}^{T},
\qquad
A_{M\times K}=X_{\mathrm{cm}}^{T},
\qquad
B_{K\times N}=A_{\mathrm{cm}}^{T}
\]

e le dimensioni passate alla libreria sono:

\[
M=k,\qquad N=\mloc,\qquad K=\nloc.
\]

Entrambi gli operandi vengono utilizzati con
\code{CUBLAS\_OP\_N}: la trasposizione non viene eseguita dalla libreria,
ma è già implicita nella reinterpretazione column-major degli stessi buffer.

\subsection{Significato del confronto}

Il backend cuBLAS costituisce un riferimento esterno e non un ulteriore
passo della sequenza di ottimizzazione sviluppata nel progetto. La libreria
sceglie internamente l'implementazione del prodotto e il progetto non
controlla direttamente parametri quali tiling, distribuzione del lavoro od
occupancy.

Il confronto viene quindi utilizzato in senso sperimentale: permette di
stabilire quale prestazione ottengano i kernel specializzati del progetto
rispetto a una routine GEMM di libreria eseguita sulla stessa GPU e con gli
stessi dati.

Nel problema considerato una delle dimensioni del prodotto, $k$, è
intenzionalmente piccola rispetto a $m$ e $n$. Questo rende il confronto
particolarmente interessante, perché i kernel sviluppati nel progetto
sono costruiti esplicitamente attorno alla struttura di un multivettore
stretto, mentre cuBLAS riceve il problema attraverso la propria interfaccia
GEMM generale. L'eventuale vantaggio di una delle due soluzioni deve però
essere stabilito dalle misure, non assunto a priori.


\section{Più rank MPI sulla stessa GPU}
\label{sec:serializzazione}

La piattaforma utilizzata per le misure dispone di una sola GPU e i backend
CUDA del progetto selezionano esplicitamente il device~0. Quando
l'applicazione viene eseguita con più rank MPI sullo stesso nodo, ogni
processo crea quindi il proprio contesto CUDA sulla medesima scheda.

In assenza di NVIDIA Multi-Process Service (MPS), il lavoro proveniente da
contesti CUDA appartenenti a processi differenti non viene eseguito
contemporaneamente sugli SM, ma viene alternato mediante time slicing.
I rank possono quindi procedere contemporaneamente sul lato CPU, ma
competono per l'esecuzione sulla stessa GPU.

Questa situazione modifica il significato delle misure ottenute mediante
\code{cudaEvent}. Per ciascun rank, $t_{\mathrm{kernel}}$ misura
l'intervallo compreso fra gli eventi che delimitano il proprio kernel.
Se il contesto viene sospeso durante tale intervallo per consentire
l'esecuzione di un altro rank, il tempo osservato risente anche della
contesa per la GPU.

Di conseguenza, con più rank per GPU, $t_{\mathrm{kernel}}$ non deve essere
interpretato come il tempo che lo stesso kernel otterrebbe disponendo della
scheda in modo esclusivo.

Il driver conserva due prospettive differenti sulla fase locale.

La prima è il tempo del kernel restituito dal backend CUDA. Al termine
delle ripetizioni, come per le altre fasi, viene considerato il massimo
fra i rank, poiché l'operazione distribuita non può essere considerata
conclusa prima del processo più lento. Questa misura rimane utile per
osservare il comportamento dei singoli backend, ma con più contesti sulla
stessa GPU è influenzata dalla loro contesa.

Per misurare invece il tempo effettivamente richiesto dall'intero gruppo
di rank per completare la fase locale, il benchmark mette a disposizione
la modalità \code{--group-timing}. In questo caso una barriera MPI viene
posta immediatamente prima di \code{local\_gemm} e una seconda subito dopo:

\[
\text{barriera}
\;\longrightarrow\;
\text{fase locale dei rank}
\;\longrightarrow\;
\text{barriera}.
\]

L'intervallo fra le due barriere misura quindi il tempo trascorso dal
momento in cui tutti i rank sono pronti fino al completamento dell'ultimo.
È questo il denominatore utilizzato per la metrica
\code{gflops\_local\_group}, che descrive il throughput complessivo del
gruppo durante la fase locale.

Per le campagne dedicate alla caratterizzazione del kernel viene quindi
utilizzato un solo rank MPI. In questo modo il backend dispone
esclusivamente della GPU e differenze dovute, per esempio, alla dimensione
del blocco CUDA, al padding della shared memory o alla scelta del layout di
$X$ non vengono confuse con la schedulazione fra contesti differenti.

A parità di dimensioni globali $M$, $N$ e $k$, il numero complessivo di
operazioni rimane lo stesso, ma con un solo rank cambia naturalmente la
dimensione del blocco locale assegnato al processo. Le misure a rank
singolo devono quindi essere interpretate come caratterizzazione del
backend sulla GPU dedicata, mentre quelle multi-rank descrivono
l'esecuzione della configurazione MPI effettiva sulla piattaforma
disponibile.